\section{Modelagem do Problema}%
\label{sec:modelagem}

A Árvore Geradora Mínima (AGM) pode ser formulada por um \textit{single-commodity flow model}, no qual o fluxo enviado pelo nó raiz garante conectividade e impede ciclos sem usar restrições de corte.

Seja $G = (V,E)$ um grafo não direcionado, onde $V = \{0,1,\ldots,n-1\}$ e o vértice $0$ atua como raiz.

\subsection{Variáveis de Decisão}

\begin{itemize}
    \item $x_{uv} \in \{0,1\}$: indica se a aresta $(u,v)$ faz parte da AGM.\@
    \item $f_{uv} \ge 0$: fluxo enviado de $u$ para $v$, permitido apenas em arestas ativas.
\end{itemize}

\subsection{Função Objetivo}

\begin{equation}
    \min \sum_{(u,v)\in E} w_{uv}\, x_{uv}
\end{equation}

\subsection{Restrições}

\paragraph{(a) Conservação de Fluxo}

Na raiz:

\begin{equation}
    \sum_{v \in V} f_{0v} - \sum_{v \in V} f_{v0} = n-1
\end{equation}

Nos demais vértices:

\begin{equation}
    \sum_{u \in V} f_{u v} - \sum_{w \in V} f_{v w} = 1,
    \quad \forall v \in V \setminus \{0\}
\end{equation}

\paragraph{(b) Capacidades (ligação fluxo–aresta)}

\begin{equation}
    f_{uv} \le (n-1)x_{uv}, \quad \forall (u,v)\in E
\end{equation}

\begin{equation}
    f_{vu} \le (n-1)x_{uv}, \quad \forall (u,v)\in E
\end{equation}

\paragraph{(c) Cardinalidade das Arestas}

\begin{equation}
    \sum_{(u,v)\in E} x_{uv} = n-1
\end{equation}

\paragraph{(d) Domínio das Variáveis}

\begin{align}
    x_{uv} &\in \{0,1\}, && \forall (u,v)\in E \\
    f_{uv} &\ge 0,       && \forall (u,v)\in E
\end{align}

\noindent
Essa formulação garante que a solução seja uma árvore conexa, acíclica e de custo mínimo.
